% !TeX root = ../main.tex

\documentclass[../main.tex]{subfiles}
\begin{document}

\chapter{ĐÁNH GIÁ TÁC ĐỘNG MÔI TRƯỜNG VÀ HIỆU QUẢ DỰ ÁN}

\section{Đánh giá và kiểm soát tác động môi trường}

Dự án nhà máy xử lý rác thải sinh hoạt phát điện mang lại lợi ích rất lớn về mặt làm sạch không gian đô thị, nhưng bản thân quá trình vận hành nhà máy cũng phát sinh ra các nguồn ô nhiễm thứ cấp (khí thải, nước thải, tro xỉ). Việc đánh giá chính xác và thiết lập hệ thống kiểm soát khắt khe là yếu tố tiên quyết để đảm bảo tính bền vững của dự án.

\subsection{Kiểm soát chất thải rắn và tro xỉ}

Nguồn gốc chất thải rắn phát sinh bao gồm: phần rác vô cơ không thể xử lý nhiệt (gạch, đá, thủy tinh, xà bần) được tách ra ở công đoạn phân loại đầu vào; than sinh học (biochar) xả ra từ lò nhiệt phân; và tro bay thu thập được từ hệ thống lọc bụi tĩnh điện (ESP). 

Phần rác không cháy chiếm tỷ lệ khoảng 13\% (tương đương 16 tấn/ngày) là rác thải trơ thông thường, không chứa thành phần nguy hại. Nguồn thải này sẽ được thu gom hàng ngày bằng xe tải chuyên dụng và vận chuyển đến khu vực chôn lấp hợp vệ sinh theo đúng quy hoạch của địa phương. Đối với than biochar (khoảng 11,5 tấn/ngày), đây là loại vật liệu giàu carbon, xốp và có khả năng giữ ẩm cao. Nhà máy sẽ tiến hành phối hợp với các cơ quan kiểm định phân tích mẫu định kỳ. Nếu hàm lượng các kim loại nặng và hóa chất vô cơ đáp ứng các tiêu chuẩn quy định, biochar sẽ được đóng bao xuất bán làm nguyên liệu phân bón hữu cơ hoặc chất cải tạo đất.

Đặc biệt lưu ý là lượng tro bay từ hệ thống xử lý khí thải, do đặc thù có khả năng hấp thụ các hợp chất hữu cơ độc hại và kim loại nặng bay hơi, nên được phân loại vào nhóm chất thải có nguy cơ nguy hại. Để kiểm soát tuyệt đối, tro bay sẽ được thu gom tự động bằng hệ thống đường ống kín, đưa vào hệ thống trộn với tỷ lệ xi măng và phụ gia đóng rắn. Khối rắn sau khi đông kết sẽ ngăn chặn triệt để sự phát tán bụi và rửa trôi kim loại, sau đó được chuyển giao cho các đơn vị có chức năng xử lý chất thải nguy hại (như Công ty Cổ phần Môi trường Đô thị Quảng Nam) vận chuyển và xử lý.

\subsection{Hệ thống xử lý nước thải}

Nước thải phát sinh từ nhà máy bao gồm nước rỉ rác sinh ra từ khu vực bunker chứa rác, nước thải từ công đoạn ép tách ẩm rác cơ học, nước ngưng từ quá trình sấy, nước thải sinh hoạt của cán bộ công nhân viên và nước rửa xe, vệ sinh phân xưởng. Trong đó, nước rỉ rác là thành phần ô nhiễm nghiêm trọng nhất, chứa nồng độ COD, BOD, N-NH3 và Coliform rất cao, mang theo mùi hôi thối đặc trưng.

Giải pháp công nghệ áp dụng cho trạm xử lý nước thải tập trung là quy trình hóa lý kết hợp sinh học màng (MBR). Dòng nước rỉ rác đầu vào trước tiên được đi qua song chắn rác và bể lắng cát, sau đó đưa vào bể điều hòa để ổn định lưu lượng và nồng độ. Tiếp đó, nước được bơm lên tháp Stripping để loại bỏ phần lớn Amoni ($NH_3$). Ở công đoạn sinh học, vi sinh vật hiếu khí và thiếu khí sẽ phân hủy các chất hữu cơ hòa tan (BOD, COD). Lớp màng lọc sinh học MBR với kích thước lỗ màng siêu nhỏ sẽ giữ lại toàn bộ bùn hoạt tính và vi khuẩn, cho ra dòng nước trong suốt. Tại bước cuối cùng, nước đi qua cột lọc than hoạt tính và bể khử trùng bằng Clo. 

Nước xả sau cùng cam kết đạt tiêu chuẩn QCVN 40:2011/BTNMT (Cột A) trước khi tái sử dụng cho các mục đích làm mát thiết bị, tưới cây, rửa đường trong khuôn viên nhà máy, đảm bảo nguyên tắc "không xả thải nước công nghiệp ra môi trường ngoài" (Zero Liquid Discharge).

\subsection{Xử lý khí thải và mùi hôi}

Khí thải phát sinh chủ yếu từ cụm lò sấy nhiệt dư, ống khói lò nhiệt phân và khí xả của động cơ diesel phát điện. Mặc dù rác đã được ép viên RDF có đặc tính cháy tốt, nhưng khí thải vẫn tiềm ẩn bụi lơ lửng, khí CO, NOx, SO2 và các hợp chất Dioxin/Furan nếu nhiệt độ cháy không được kiểm soát tốt.

Hệ thống xử lý khí thải được thiết kế liên hoàn nhiều bậc. Đầu tiên, nhiệt độ dòng khí được dập nhanh (quenching) xuống dưới 200°C trong thời gian chưa tới 1 giây nhằm ngăn ngừa quá trình tái tổ hợp tạo thành Dioxin/Furan. Sau đó, khí được dẫn qua thiết bị lọc bụi tĩnh điện (ESP) với hiệu suất thu hồi hạt vi mô lên tới 99,5\%. Dòng khí sạch bụi tiếp tục được đẩy qua tháp hấp thụ kiểu ướt (Scrubber). Tại đây, dung dịch hóa chất kiềm (NaOH) được phun mù dạng sương ngược chiều với dòng khí bay lên, xảy ra phản ứng trung hòa hoàn toàn các khí mang tính axit (SO2, HCl, HF). Hệ thống xả khí tự động qua ống khói cao, có trang bị cảm biến quan trắc liên tục, đảm bảo thông số đầu ra luôn đạt tiêu chuẩn QCVN 61-MT:2016/BTNMT.

Vấn đề mùi hôi tại khu vực tiếp nhận rác được giải quyết triệt để nhờ hệ thống tạo áp suất âm liên tục trong nhà xưởng, không khí mang mùi được hút trực tiếp đưa vào buồng đốt lò làm khí cấp 1, qua đó thiêu hủy hoàn toàn các phân tử gây mùi dưới tác dụng của nhiệt độ cao.

\section{Công tác quản lý, giám sát và phòng ngừa rủi ro}

An toàn vận hành và ứng phó sự cố là ưu tiên hàng đầu tại nhà máy. Hệ thống giám sát (SCADA) được trang bị tại trung tâm điều khiển, thu thập dữ liệu tự động 24/24h từ các cảm biến nhiệt độ, áp suất, lưu lượng được phân bổ tại mọi mắt xích công nghệ. Khi có bất kỳ thông số nào vượt ngưỡng cảnh báo, hệ thống sẽ tự động kích hoạt chuông báo động, cắt nguồn cấp nhiên liệu và mở các van xả áp an toàn tĩnh điện.

Các kịch bản rủi ro nghiêm trọng như mất điện lưới đột ngột, hỏa hoạn tại kho lưu trữ RDF, hay rò rỉ khí syngas đều được lập phương án ứng phó chi tiết. Nhà máy duy trì một máy phát điện dự phòng khẩn cấp để duy trì hoạt động của bơm nước dập lửa và quạt thông gió trong mọi tình huống. Đội ngũ kỹ sư vận hành yêu cầu phải trải qua các khóa huấn luyện nghiêm ngặt về an toàn lao động, phòng cháy chữa cháy và được cấp chứng chỉ trước khi trực tiếp làm việc tại xưởng.

\section{Đánh giá hiệu quả kinh tế - xã hội}

\subsection{Hiệu quả kinh tế}

Việc chuyển đổi sang công nghệ nhiệt phân kết hợp động cơ diesel thể hiện hiệu quả kinh tế rõ rệt thông qua việc tối ưu hóa chi phí đầu tư và chi phí vận hành (OPEX). So với công nghệ khí hóa và tuabin khí đòi hỏi chi phí đầu tư thiết bị tinh sạch cao và phí bảo trì đắt đỏ, mô-đun nhiệt phân và động cơ đốt trong thân thiện hơn về mặt tài chính và phù hợp với nguồn nhân lực hiện tại. 

Nguồn thu của dự án khá đa dạng và ổn định: ngoài nguồn thu từ phí xử lý rác do ngân sách địa phương chi trả theo quy định, nhà máy còn có nguồn doanh thu đáng kể từ việc bán sản lượng điện 1,82 MW lên lưới quốc gia, và nguồn thu phụ từ việc bán than biochar cho ngành nông nghiệp cũng như bán phế liệu kim loại. Các luồng tiền này giúp nhà máy đẩy nhanh thời gian hoàn vốn đầu tư và đảm bảo lợi nhuận tự duy trì lâu dài.

\subsection{Hiệu quả xã hội và môi trường}

Dự án trực tiếp giải bài toán bức thiết về rác thải sinh hoạt ùn ứ của thành phố du lịch Hội An. Nhờ có nhà máy, địa phương chấm dứt tình trạng phải vận chuyển rác đi hơn 100 km vào bãi chôn lấp Tam Nghĩa, tiết kiệm một lượng ngân sách vận tải khổng lồ và chấm dứt nguy cơ rò rỉ rác trên các tuyến quốc lộ.

Việc vận hành nhà máy tạo ra công ăn việc làm trực tiếp cho hàng trăm lao động địa phương (từ công nhân phân loại, kỹ sư vận hành đến bộ phận hành chính). Đồng thời, thông qua việc đóng góp vào quỹ đất sạch (không phải dùng quỹ đất làm bãi chôn lấp), dự án kiến tạo môi trường trong lành, giữ gìn hình ảnh mỹ quan đô thị sạch đẹp cho Hội An, qua đó góp phần thúc đẩy ngành dịch vụ du lịch phát triển mạnh mẽ và bền vững hơn.

\end{document}